\section{Introduction to the Control Plane (5.1)}
%=============================================================================

\subsection{Network-Layer Functions Recap}

\begin{keybox}[Forwarding vs.\ Routing -- Two Key Functions]
The network layer has two fundamental functions:
\begin{itemize}[noitemsep]
    \item \textbf{Forwarding:} Move packets from a router's input port to the appropriate output port. This is a \emph{local}, per-router action belonging to the \textbf{data plane}.
    \item \textbf{Routing:} Determine the route (path) that packets take from source to destination. This is a \emph{network-wide}, global action belonging to the \textbf{control plane}.
\end{itemize}

\medskip
\textbf{Analogy:} Forwarding = taking the correct exit at a roundabout. Routing = planning the entire trip from city A to city B.
\end{keybox}

\subsection{Two Approaches to the Control Plane}

\begin{keybox}[Per-Router Control vs.\ SDN]
There are two main approaches to structuring the network control plane:

\textbf{1. Per-Router Control (Traditional):}
\begin{itemize}[noitemsep]
    \item Individual routing algorithm components run in \emph{each and every router}
    \item Routers interact with each other in the control plane to compute forwarding tables
    \item Each router independently runs a routing algorithm (e.g., OSPF, BGP)
\end{itemize}

\textbf{2. Logically Centralized Control (SDN):}
\begin{itemize}[noitemsep]
    \item A remote controller computes and distributes forwarding tables to routers
    \item Routers/switches receive their flow tables from the central controller
    \item Physically separates the control plane from the data plane
\end{itemize}
\end{keybox}

\begin{imagebox}[Per-Router Control vs.\ SDN Control Plane]
\figureimgpair{Figure 5.1.png}{Figure 5.2.png}
\end{imagebox}

%=============================================================================
\section{Routing Algorithms (5.2)}
%=============================================================================

\subsection{Goal and Overview}

\begin{keybox}[Routing Protocol Goal]
The goal of a routing protocol is to determine ``good'' paths (routes) from sending hosts to receiving hosts through the network of routers.
\begin{itemize}[noitemsep]
    \item A \textbf{path} is a sequence of routers that packets traverse from source to destination
    \item ``Good'' typically means: least cost, fastest, or least congested
    \item Routing is considered a ``top-10'' networking challenge
\end{itemize}
\end{keybox}

\subsection{Graph Abstraction}

\begin{keybox}[Network as a Graph]
A network can be modelled as a graph $G = (N, E)$:
\begin{itemize}[noitemsep]
    \item $N$ = set of routers (nodes): $\{u, v, w, x, y, z\}$
    \item $E$ = set of links (edges): $\{(u,v), (u,x), (v,x), \ldots\}$
    \item Each edge has a \textbf{cost} $c(x,y)$ (e.g., bandwidth, delay, monetary cost, congestion)
    \item $c(x,y) = \infty$ if $(x,y)$ is not a direct link
    \item The cost of a path $(x_1, x_2, \ldots, x_p) = c(x_1,x_2) + c(x_2,x_3) + \cdots + c(x_{p-1},x_p)$
\end{itemize}

\textbf{Key question:} What is the least-cost path between two nodes?
\end{keybox}

\subsection{Routing Algorithm Classification}

\begin{exambox}[Exam-relevant: Classification of Routing Algorithms]
Routing algorithms can be classified along two axes:

\textbf{Axis 1 -- Information scope:}
\begin{itemize}[noitemsep]
    \item \textbf{Global (link state):} All routers have \emph{complete} topology and link cost information. Each router computes shortest paths independently using global knowledge.
    \item \textbf{Decentralized (distance vector):} Iterative process of computation and exchange of information with neighbors. Routers initially only know link costs to their directly attached neighbors.
\end{itemize}

\textbf{Axis 2 -- How fast do routes change?}
\begin{itemize}[noitemsep]
    \item \textbf{Static:} Routes change slowly over time (manual configuration)
    \item \textbf{Dynamic:} Routes change more quickly -- periodic updates or in response to link cost changes
\end{itemize}
\end{exambox}

%=============================================================================
\section{Link-State Routing Algorithm: Dijkstra's Algorithm (5.2.1)}
%=============================================================================

\subsection{Overview}

\begin{keybox}[Link-State Algorithm -- Key Properties]
In a \textbf{link-state} algorithm:
\begin{itemize}[noitemsep]
    \item All routers have \emph{complete} knowledge of the network topology and all link costs
    \item This is achieved by \textbf{link-state broadcast} -- each router broadcasts its link information to \emph{all} other routers in the network (flooding)
    \item Each router then independently runs \textbf{Dijkstra's algorithm} to compute shortest paths from itself to all other nodes
    \item The result is a \textbf{shortest-path tree} rooted at the computing router, which determines the forwarding table
\end{itemize}
\end{keybox}

\subsection{Dijkstra's Algorithm}

\begin{formulabox}[Dijkstra's Algorithm -- Notation and Procedure]
\textbf{Notation:}
\begin{itemize}[noitemsep]
    \item $c(x,y)$: direct link cost from node $x$ to node $y$; $= \infty$ if not direct neighbors
    \item $D(v)$: current estimate of the least-cost path from source to $v$
    \item $p(v)$: predecessor node along the current least-cost path to $v$
    \item $N'$: set of nodes whose least-cost path is definitively known
\end{itemize}

\textbf{Algorithm:}
\begin{enumerate}[noitemsep]
    \item \textbf{Initialize:} $N' = \{u\}$ (source node). For all neighbors $v$ of $u$: $D(v) = c(u,v)$. For all non-neighbors: $D(v) = \infty$.
    \item \textbf{Loop:} Find node $w \notin N'$ such that $D(w)$ is minimum. Add $w$ to $N'$.
    \item \textbf{Update:} For each neighbor $v$ of $w$ that is not in $N'$:
    \[
    D(v) = \min\big(D(v),\; D(w) + c(w,v)\big)
    \]
    If the path through $w$ is better, update $p(v) = w$.
    \item \textbf{Repeat} step 2--3 until all nodes are in $N'$.
\end{enumerate}

\textbf{Complexity:} $O(n^2)$ for $n$ nodes (can be improved to $O(n \log n)$ with a priority queue).
\end{formulabox}

\begin{exambox}[Exam-relevant: Dijkstra's Algorithm Walkthrough]
\textbf{Example:} Given the following network (from textbook Figure 5.3):

\begin{center}
\begin{tabular}{cccccccc}
\toprule
\textbf{Step} & $N'$ & $D(v),p(v)$ & $D(w),p(w)$ & $D(x),p(x)$ & $D(y),p(y)$ & $D(z),p(z)$ \\
\midrule
0 & $\{u\}$ & $2,u$ & $5,u$ & $1,u$ & $\infty$ & $\infty$ \\
1 & $\{u,x\}$ & $2,u$ & $4,x$ & -- & $2,x$ & $\infty$ \\
2 & $\{u,x,v\}$ & -- & $4,x$ & -- & $2,x$ & $\infty$ \\
3 & $\{u,x,v,y\}$ & -- & $4,x$ & -- & -- & $4,y$ \\
4 & $\{u,x,v,y,w\}$ & -- & -- & -- & -- & $4,y$ \\
5 & $\{u,x,v,y,w,z\}$ & -- & -- & -- & -- & -- \\
\bottomrule
\end{tabular}
\end{center}

At each step, the node with the smallest $D(\cdot)$ not yet in $N'$ is added, and all neighbor distances are updated.

\medskip
\textbf{Tip:} On the exam, you may be asked to run Dijkstra's algorithm on a given graph. Practice filling out a table like the one above step by step.
\end{exambox}

\begin{imagebox}[Dijkstra's Algorithm Example]
\figureimgpair{Figure 5.3.png}{Figure 5.4.png}
\end{imagebox}

\subsection{Oscillations in Link-State Routing}

\begin{warnbox}[Oscillation Problem]
When link costs are \textbf{traffic-dependent} (e.g., proportional to the traffic load on a link), Dijkstra's algorithm can cause \textbf{oscillations}:
\begin{itemize}[noitemsep]
    \item All routers compute new routes based on current traffic
    \item Traffic shifts to the new ``best'' path, making it congested
    \item In the next iteration, the old path becomes least-cost again
    \item Routes keep oscillating back and forth
\end{itemize}

\textbf{Mitigation:} Routers can randomize the timing of their link-state broadcasts so they don't all recompute routes simultaneously.
\end{warnbox}

%=============================================================================
\section{Distance Vector Routing Algorithm (5.2.2)}
%=============================================================================

\subsection{Bellman-Ford Equation}

\begin{formulabox}[Bellman-Ford Equation (Dynamic Programming)]
The Bellman-Ford equation defines the least-cost path recursively:
\[
d_x(y) = \min_v \big\{ c(x,v) + d_v(y) \big\}
\]
where:
\begin{itemize}[noitemsep]
    \item $d_x(y)$: cost of the least-cost path from $x$ to $y$
    \item $\min_v$: minimum taken over all neighbors $v$ of $x$
    \item $c(x,v)$: cost of the direct link from $x$ to neighbor $v$
    \item $d_v(y)$: cost of the least-cost path from $v$ to $y$ (as reported by neighbor $v$)
\end{itemize}

The neighbor $v^*$ that achieves the minimum becomes the \textbf{next hop} in the forwarding table for destination $y$.
\end{formulabox}

\subsection{Distance Vector Algorithm}

\begin{keybox}[DV Algorithm -- How It Works]
The distance vector algorithm is:
\begin{itemize}[noitemsep]
    \item \textbf{Iterative:} Continues until no more information is exchanged
    \item \textbf{Asynchronous:} Nodes do not need to operate in lockstep
    \item \textbf{Distributed:} Each node only communicates with its direct neighbors
\end{itemize}

\textbf{Process:}
\begin{enumerate}[noitemsep]
    \item Each node $x$ maintains a \textbf{distance vector} $\mathbf{D}_x = [D_x(y) : y \in N]$ -- its estimate of the least cost to every other node
    \item Each node $x$ maintains its neighbors' distance vectors: $\mathbf{D}_v$ for each neighbor $v$
    \item Periodically, or when its local DV changes, each node sends its DV to its neighbors
    \item When node $x$ receives a new DV from a neighbor $v$, it updates its own DV using Bellman-Ford:
    \[
    D_x(y) = \min_v \big\{ c(x,v) + D_v(y) \big\} \quad \text{for each node } y
    \]
    \item If $D_x(y)$ changes, $x$ notifies its neighbors
\end{enumerate}

The algorithm converges to the correct least-cost paths. Each node only needs to know its own link costs and its neighbors' distance vectors.
\end{keybox}

\begin{imagebox}[Distance-Vector Algorithm]
\figureimg{Figure 5.6.png}
\end{imagebox}

\subsection{Count-to-Infinity Problem}

\begin{exambox}[Exam-relevant: Count-to-Infinity Problem]
A major issue with distance vector routing is the \textbf{count-to-infinity problem}:

\textbf{Scenario:} Suppose a link cost increases dramatically (e.g., $y$--$x$ cost goes from 4 to 60):
\begin{itemize}[noitemsep]
    \item Node $y$ sees the cost to $x$ via direct link is now 60
    \item But $y$ thinks it can reach $x$ via $z$ for cost $5$ (because $z$'s old DV says $D_z(x) = 5$)
    \item So $y$ sets $D_y(x) = c(y,z) + D_z(x) = 1 + 5 = 6$
    \item When $z$ gets $y$'s update, $z$ updates $D_z(x) = c(z,y) + D_y(x) = 1 + 6 = 7$
    \item This continues: $D_y(x) = 8$, $D_z(x) = 9$, $D_y(x) = 10$, \ldots
    \item It takes \textbf{44 iterations} to converge to the correct value of 60!
\end{itemize}

This is called ``bad news travels slowly.'' Good news (a cost decrease) propagates in one iteration.

\medskip
\textbf{Mitigation -- Poisoned Reverse:}
\begin{itemize}[noitemsep]
    \item If $z$ routes to $x$ through $y$, then $z$ tells $y$ that $D_z(x) = \infty$
    \item This prevents $y$ from routing to $x$ through $z$ (avoiding the loop)
    \item \textbf{Limitation:} Poisoned reverse only solves loops involving two nodes. Loops involving three or more nodes are \emph{not} detected.
\end{itemize}
\end{exambox}

%=============================================================================
\section{Comparison: Link State vs.\ Distance Vector (5.2.3)}
%=============================================================================

\begin{exambox}[Exam-relevant: LS vs.\ DV Comparison]
\begin{center}
\begin{tabular}{p{3.5cm}p{5.5cm}p{5.5cm}}
\toprule
\textbf{Property} & \textbf{Link State (LS)} & \textbf{Distance Vector (DV)} \\
\midrule
\textbf{Message complexity} & $O(n \cdot E)$ messages sent (each of $n$ nodes broadcasts to $E$ neighbors) & Exchange only between neighbors; convergence time varies \\
\midrule
\textbf{Speed of convergence} & $O(n^2)$ algorithm; may have oscillations with traffic-dependent costs & Convergence time varies; count-to-infinity problem; routing loops possible \\
\midrule
\textbf{Robustness (what if a router malfunctions?)} & A router can advertise \emph{incorrect link cost}; each node computes its own table -- errors are somewhat contained & A router can advertise \emph{incorrect path cost}; errors propagate through the network (``rumors'') \\
\bottomrule
\end{tabular}
\end{center}

\medskip
\textbf{Key takeaway:}
\begin{itemize}[noitemsep]
    \item LS: each node computes independently based on global information -- errors are localized
    \item DV: nodes rely on neighbors' information -- a single incorrect value can propagate everywhere
\end{itemize}
\end{exambox}

%=============================================================================
\section{Intra-AS Routing: OSPF (5.3)}
%=============================================================================

\subsection{Autonomous Systems and Scalability}

\begin{keybox}[Why Hierarchical Routing?]
A flat routing scheme where all routers run the same algorithm does not scale because:
\begin{itemize}[noitemsep]
    \item \textbf{Scale:} Billions of destinations -- storing all destinations in routing tables is infeasible; routing message exchange would overwhelm links
    \item \textbf{Administrative autonomy:} Each network admin wants to control routing within their own network
\end{itemize}

\textbf{Solution: Autonomous Systems (AS)}
\begin{itemize}[noitemsep]
    \item The Internet is divided into \textbf{autonomous systems} (ASes)
    \item An AS is a collection of routers under the \emph{same administrative control} (e.g., an ISP, a university, a company)
    \item \textbf{Intra-AS routing:} Routing \emph{within} an AS (all routers run the same intra-AS protocol)
    \item \textbf{Inter-AS routing:} Routing \emph{between} ASes (handled by a separate protocol -- BGP)
    \item \textbf{Gateway routers:} Routers at the ``edge'' of an AS that connect to routers in other ASes
\end{itemize}
\end{keybox}

\subsection{OSPF -- Open Shortest Path First}

\begin{keybox}[OSPF -- Overview]
OSPF is the most widely used \textbf{intra-AS} (interior gateway) routing protocol. It uses a \textbf{link-state} algorithm.

\textbf{Key characteristics:}
\begin{itemize}[noitemsep]
    \item Uses \textbf{Dijkstra's algorithm} to compute shortest-path tree
    \item Each router floods \textbf{link-state advertisements (LSAs)} to \emph{all} other routers in the AS
    \item LSAs contain: one entry per directly attached neighbor, with the link cost
    \item OSPF messages are carried directly over \textbf{IP} (protocol number 89) -- not TCP or UDP
    \item ``Open'' means the protocol specification is publicly available
\end{itemize}

\textbf{Advanced features of OSPF:}
\begin{itemize}[noitemsep]
    \item \textbf{Security:} All OSPF messages can be authenticated (prevents malicious intrusion)
    \item \textbf{Multiple same-cost paths:} When multiple paths have equal cost, OSPF allows traffic to be spread across them (ECMP -- Equal-Cost Multi-Path)
    \item \textbf{Integrated unicast and multicast:} Multicast OSPF (MOSPF) uses the same topology database
    \item \textbf{Hierarchical OSPF:} Support for large ASes through \emph{areas}
\end{itemize}
\end{keybox}

\subsection{Hierarchical OSPF}

\begin{keybox}[Two-Level Hierarchy in OSPF]
For large ASes, OSPF supports a two-level hierarchy:
\begin{itemize}[noitemsep]
    \item \textbf{Local areas:} Each area runs its own OSPF link-state algorithm. Routers only flood LSAs \emph{within} their area.
    \item \textbf{Backbone area (Area 0):} Connects all local areas. Used to route traffic \emph{between} areas.
    \item \textbf{Area border routers:} Belong to both a local area and the backbone. They summarize routing information from their area and advertise it to the backbone.
    \item \textbf{Backbone routers:} Run OSPF routing limited to the backbone area.
    \item \textbf{Boundary routers:} Connect to other ASes (run inter-AS routing as well).
\end{itemize}
\end{keybox}

\begin{imagebox}[OSPF Autonomous System with Areas]
\end{imagebox}

%=============================================================================
\section{Inter-AS Routing: BGP (5.4)}
%=============================================================================

\subsection{BGP Overview}

\begin{keybox}[BGP -- Border Gateway Protocol]
BGP is the \textbf{inter-AS} (inter-domain) routing protocol -- the ``glue that holds the Internet together.''

\textbf{BGP provides each AS with a means to:}
\begin{itemize}[noitemsep]
    \item \textbf{eBGP} (external BGP): Obtain subnet reachability information from neighboring ASes
    \item \textbf{iBGP} (internal BGP): Propagate reachability information to all routers \emph{within} the AS
    \item Determine ``good'' routes to other networks based on reachability information and \textbf{policy}
\end{itemize}

\textbf{Key facts:}
\begin{itemize}[noitemsep]
    \item BGP allows each subnet to advertise its existence to the rest of the Internet: ``I am here, and here is how to reach me''
    \item BGP runs over \textbf{TCP} (port 179) -- reliable message exchange
    \item BGP sessions are called \textbf{BGP connections} (semi-permanent TCP connections)
\end{itemize}
\end{keybox}

\subsection{BGP Basics: Path Attributes and Routes}

\begin{keybox}[BGP Route Advertisements]
When a router advertises a prefix via BGP, it includes \textbf{BGP attributes}. A prefix together with its attributes is called a \textbf{route}.

\textbf{Two important attributes:}
\begin{itemize}[noitemsep]
    \item \textbf{AS-PATH:} The list of ASes through which the advertisement has passed. Example: AS2 AS17 AS67. Used for loop detection -- if a router sees its own AS in the path, it rejects the advertisement.
    \item \textbf{NEXT-HOP:} The IP address of the router interface that begins the AS-PATH. Indicates the next-hop router to reach the destination.
\end{itemize}

\textbf{Policy-based routing:}
\begin{itemize}[noitemsep]
    \item A gateway receiving a route advertisement uses its \textbf{import policy} to decide whether to accept or filter the route
    \item An AS uses its \textbf{export policy} to decide whether to advertise routes to neighboring ASes
\end{itemize}
\end{keybox}

\begin{imagebox}[eBGP and iBGP Connections]
\figureimg{Figure 5.9.png}
\end{imagebox}

\subsection{BGP Route Selection}

\begin{exambox}[Exam-relevant: BGP Route Selection Rules]
When a router learns about multiple routes to the same prefix, it selects the best route using the following criteria (in order of priority):

\begin{enumerate}[noitemsep]
    \item \textbf{Local preference value} (policy decision): Highest local preference wins. Set by the AS administrator.
    \item \textbf{Shortest AS-PATH:} Fewest number of ASes traversed.
    \item \textbf{Closest NEXT-HOP router} (hot potato routing): Choose the route whose NEXT-HOP router has the lowest intra-AS cost from the current router.
    \item \textbf{Additional criteria:} BGP identifiers and other tie-breakers.
\end{enumerate}

\medskip
\textbf{Hot Potato Routing:} The router chooses the gateway that has the \emph{smallest least cost} (determined by the intra-AS algorithm) to reach, regardless of the inter-AS cost. The idea: ``get rid of the packet as quickly as possible'' (within the AS).
\end{exambox}

\subsection{Why Different Intra-AS and Inter-AS Routing?}

\begin{keybox}[Intra-AS vs.\ Inter-AS: Different Goals]
\textbf{Policy:}
\begin{itemize}[noitemsep]
    \item \textbf{Inter-AS (BGP):} Policy is dominant. An admin wants control over how traffic is routed and who routes through their network. BGP carries path attributes and allows filtering.
    \item \textbf{Intra-AS (OSPF):} Single admin, so policy is less of an issue. Focus on performance.
\end{itemize}

\textbf{Scale:}
\begin{itemize}[noitemsep]
    \item Hierarchical routing saves table size and reduces update traffic
    \item Intra-AS can focus on performance; inter-AS may sacrifice performance for policy
\end{itemize}

\textbf{Performance:}
\begin{itemize}[noitemsep]
    \item Intra-AS routing can focus on finding optimal paths
    \item Inter-AS routing is constrained by policies (may not use shortest path)
\end{itemize}
\end{keybox}

%=============================================================================
\section{ICMP -- Internet Control Message Protocol (5.6)}
%=============================================================================

\subsection{ICMP Overview}

\begin{keybox}[ICMP -- Purpose and Position]
ICMP is used by hosts and routers to communicate \textbf{network-level information}:
\begin{itemize}[noitemsep]
    \item \textbf{Error reporting:} Unreachable host, network, port, or protocol
    \item \textbf{Echo request/reply:} Used by \texttt{ping}
    \item ICMP is a network-layer protocol but sits ``above'' IP:
    \begin{itemize}[noitemsep]
        \item ICMP messages are carried as the \textbf{payload} of IP datagrams
    \end{itemize}
    \item Each ICMP message has a \textbf{type} and a \textbf{code} field, plus the first 8 bytes of the IP datagram that caused the error
\end{itemize}
\end{keybox}

\subsection{ICMP Message Types}

\begin{exambox}[Exam-relevant: ICMP Message Types and Codes]
\begin{center}
\begin{tabular}{ccl}
\toprule
\textbf{Type} & \textbf{Code} & \textbf{Description} \\
\midrule
0 & 0 & Echo reply (ping response) \\
3 & 0 & Destination network unreachable \\
3 & 1 & Destination host unreachable \\
3 & 2 & Destination protocol unreachable \\
3 & 3 & Destination port unreachable \\
3 & 6 & Destination network unknown \\
3 & 7 & Destination host unknown \\
4 & 0 & Source quench (congestion control -- not used) \\
8 & 0 & Echo request (ping) \\
9 & 0 & Route advertisement \\
10 & 0 & Router discovery \\
11 & 0 & TTL expired \\
12 & 0 & Bad IP header \\
\bottomrule
\end{tabular}
\end{center}

\textbf{Key types to remember:}
\begin{itemize}[noitemsep]
    \item Type 8/0: Echo request and reply (\texttt{ping})
    \item Type 11: TTL expired (used by \texttt{traceroute})
    \item Type 3: Destination unreachable (various codes for network, host, port, protocol)
\end{itemize}
\end{exambox}

\begin{imagebox}[ICMP Message Types]
\figureimg{Figure 5.19.png}
\end{imagebox}

\subsection{Traceroute and ICMP}

\begin{exambox}[Exam-relevant: How Traceroute Works]
\texttt{Traceroute} uses ICMP to discover the path packets take through the network:

\textbf{How it works:}
\begin{enumerate}[noitemsep]
    \item The source sends sets of \textbf{UDP segments} to the destination
    \item The 1st set has \textbf{TTL = 1}, the 2nd set has \textbf{TTL = 2}, etc. (typically 3 probes per TTL value)
    \item When a datagram arrives at the $n$th router, its TTL expires (reaches 0)
    \item The router \textbf{discards} the datagram and sends back an \textbf{ICMP ``TTL expired'' message} (type 11, code 0) to the source
    \item The ICMP message includes the router's name and IP address
    \item When the ICMP message arrives at the source, it records the \textbf{RTT} (round-trip time)
\end{enumerate}

\textbf{Stopping criteria:}
\begin{itemize}[noitemsep]
    \item Eventually a UDP segment reaches the destination host
    \item The destination returns an ICMP \textbf{``port unreachable''} message (type 3, code 3) because the destination port is unlikely to be in use
    \item The source receives this message and knows the destination has been reached -- it stops
\end{itemize}

\medskip
\textbf{Output:} For each TTL value, traceroute displays the router IP address and three RTT measurements.
\end{exambox}

%=============================================================================
\section{SDN -- Software-Defined Networking (5.5)}
%=============================================================================

\subsection{SDN Overview}

\begin{keybox}[SDN Control Plane]
In Software-Defined Networking, the control plane is \textbf{logically centralized}:
\begin{itemize}[noitemsep]
    \item A remote \textbf{SDN controller} computes and installs forwarding tables (flow tables) in each router/switch
    \item The controller has a \textbf{global view} of the network (topology, traffic, link states)
    \item Routers/switches become simple forwarding devices (``dumb'' data plane)
    \item The controller is implemented as a distributed system for reliability and scalability
\end{itemize}

\textbf{Why SDN?}
\begin{itemize}[noitemsep]
    \item Easier network management -- one central point of configuration
    \item Table-based forwarding allows ``programming'' of routers via an open API (e.g., OpenFlow)
    \item Enables new routing algorithms without changing router hardware/software
    \item Centralized control makes it easier to implement complex policies, load balancing, and traffic engineering
\end{itemize}
\end{keybox}

\subsection{SDN Architecture}

\begin{keybox}[SDN Layers]
The SDN architecture has three main layers:
\begin{enumerate}[noitemsep]
    \item \textbf{Data plane (SDN-controlled switches):}
    \begin{itemize}[noitemsep]
        \item Fast, simple forwarding devices
        \item Flow table computed and installed by the controller
        \item API for communication with controller (e.g., OpenFlow)
    \end{itemize}
    \item \textbf{SDN Controller (network OS):}
    \begin{itemize}[noitemsep]
        \item Maintains network state information
        \item Interacts with switches below (southbound API) and applications above (northbound API)
        \item Implemented as a distributed system (e.g., ONOS, OpenDaylight)
    \end{itemize}
    \item \textbf{Network-control applications:}
    \begin{itemize}[noitemsep]
        \item ``Brain'' of the control -- routing, access control, load balancing
        \item Implemented as programs on top of the controller using the northbound API
        \item Can be provided by third parties (separate from controller and switches)
    \end{itemize}
\end{enumerate}
\end{keybox}

\begin{imagebox}[SDN Controller Architecture]
\figureimg{Figure 5.14.png}
\end{imagebox}

\subsection{OpenFlow Protocol}

\begin{keybox}[OpenFlow -- Controller-to-Switch Communication]
OpenFlow is the protocol used between the SDN controller and the switches:
\begin{itemize}[noitemsep]
    \item Operates over \textbf{TCP} (with optional encryption)
    \item \textbf{Controller-to-switch messages:} configure, modify-state (flow table entries), read-state (statistics), send-packet
    \item \textbf{Switch-to-controller messages:} packet-in (forward a packet to controller), flow-removed (flow table entry removed), port-status (port change notification)
\end{itemize}

\textbf{Generalized forwarding:} OpenFlow flow table entries have the form:
\[
\text{match} + \text{action} + \text{priority} + \text{counters}
\]
\begin{itemize}[noitemsep]
    \item \textbf{Match:} Pattern on packet header fields (source/dest IP, port, protocol, VLAN, etc.)
    \item \textbf{Action:} Forward to port(s), drop, modify fields, send to controller
    \item This generalizes traditional destination-based forwarding
\end{itemize}
\end{keybox}

%=============================================================================
\section{Network Management and SNMP (5.7)}
%=============================================================================

\subsection{Network Management Overview}

\begin{keybox}[What is Network Management?]
An autonomous system (``network'') consists of 1000s of interacting hardware/software components. Network management involves:
\begin{itemize}[noitemsep]
    \item \textbf{Deployment, integration, and coordination} of hardware, software, and human elements
    \item \textbf{Monitor, test, poll, configure, analyze, evaluate, and control} the network
    \item Goal: meet real-time operational performance and Quality of Service requirements at a reasonable cost
\end{itemize}
\end{keybox}

\subsection{SNMP -- Simple Network Management Protocol}

\begin{keybox}[SNMP]
SNMP is the protocol used for network management:
\begin{itemize}[noitemsep]
    \item \textbf{Management server} (managing entity) queries and controls managed devices
    \item \textbf{Managed devices} (routers, switches, hosts) contain a \textbf{Management Information Base (MIB)} -- a collection of objects/variables representing the device's state
    \item SNMP has two modes:
    \begin{itemize}[noitemsep]
        \item \textbf{Request/response:} Manager queries or sets MIB values on a managed device
        \item \textbf{Trap messages:} Managed device sends an unsolicited notification to the manager (e.g., link failure, threshold exceeded)
    \end{itemize}
    \item SNMP runs over \textbf{UDP}
    \item SNMP message types: GetRequest, SetRequest, GetBulkRequest, InformRequest, Response, Trap
\end{itemize}
\end{keybox}

%=============================================================================
\section{Summary and Exam Tips}
%=============================================================================

\begin{exambox}[Frequently Tested Topics -- Chapter 5]
Based on past exams and lecture emphasis:

\begin{enumerate}[noitemsep]
    \item \textbf{Dijkstra's algorithm:} Be able to run it step-by-step on a given graph and fill in the table. Know the formula: $D(v) = \min(D(v), D(w) + c(w,v))$.
    \item \textbf{Bellman-Ford equation:} $d_x(y) = \min_v \{c(x,v) + d_v(y)\}$. Understand what each term means.
    \item \textbf{LS vs.\ DV:} Know the differences in message complexity, convergence speed, and robustness.
    \item \textbf{Count-to-infinity:} Understand why it happens, how poisoned reverse helps, and its limitation.
    \item \textbf{OSPF:} Link-state protocol, runs Dijkstra's, floods LSAs, runs directly over IP (protocol 89), supports hierarchical areas.
    \item \textbf{BGP:} Inter-AS routing, uses TCP port 179, eBGP vs.\ iBGP, AS-PATH and NEXT-HOP attributes, route selection rules, hot potato routing.
    \item \textbf{ICMP:} Know the main message types (echo request/reply, TTL expired, dest unreachable). Understand how \texttt{traceroute} uses ICMP.
    \item \textbf{SDN:} Logically centralized control plane, OpenFlow, match+action paradigm, three-layer architecture.
    \item \textbf{Intra-AS vs.\ Inter-AS:} Why they are different (policy, scale, performance).
\end{enumerate}
\end{exambox}

\begin{warnbox}[Common Exam Pitfalls]
\begin{itemize}[noitemsep]
    \item \textbf{Link State $\neq$ Distance Vector} -- Don't confuse which algorithm uses global vs.\ local information. LS = global/Dijkstra, DV = local/Bellman-Ford.
    \item \textbf{OSPF runs over IP}, not TCP or UDP. BGP runs over \textbf{TCP}.
    \item \textbf{Traceroute uses UDP segments}, not ICMP echo requests. The ICMP messages (TTL expired, port unreachable) are the \emph{responses}.
    \item \textbf{Poisoned reverse does not solve all loops} -- only prevents two-node loops. Three-or-more-node loops can still occur.
    \item \textbf{BGP is policy-based}, not purely shortest-path. Local preference can override shortest AS-PATH.
    \item \textbf{Hot potato routing} minimizes \emph{intra-AS} cost to the exit point, not the total end-to-end cost.
    \item \textbf{ICMP messages are carried inside IP datagrams} -- ICMP is ``above'' IP but is considered a network-layer protocol.
    \item \textbf{eBGP vs.\ iBGP:} eBGP is between routers in \emph{different} ASes; iBGP is between routers in the \emph{same} AS. Don't mix them up.
\end{itemize}
\end{warnbox}

\begin{formulabox}[Key Formulas and Equations -- Chapter 5]
\begin{align}
\text{Dijkstra's update:} \quad D(v) &= \min\big(D(v),\; D(w) + c(w,v)\big) \\
\text{Bellman-Ford:} \quad d_x(y) &= \min_v \big\{ c(x,v) + d_v(y) \big\} \\
\text{Dijkstra complexity:} \quad &O(n^2) \text{ (or } O(n \log n) \text{ with priority queue)} \\
\text{LS message complexity:} \quad &O(n \cdot E) \text{ messages}
\end{align}
\end{formulabox}
